% Options for packages loaded elsewhere
\PassOptionsToPackage{unicode}{hyperref}
\PassOptionsToPackage{hyphens}{url}
%
\documentclass[
]{article}
\usepackage{amsmath,amssymb}
\usepackage{iftex}
\ifPDFTeX
  \usepackage[T1]{fontenc}
  \usepackage[utf8]{inputenc}
  \usepackage{textcomp} % provide euro and other symbols
\else % if luatex or xetex
  \usepackage{unicode-math} % this also loads fontspec
  \defaultfontfeatures{Scale=MatchLowercase}
  \defaultfontfeatures[\rmfamily]{Ligatures=TeX,Scale=1}
\fi
\usepackage{lmodern}
\ifPDFTeX\else
  % xetex/luatex font selection
\fi
% Use upquote if available, for straight quotes in verbatim environments
\IfFileExists{upquote.sty}{\usepackage{upquote}}{}
\IfFileExists{microtype.sty}{% use microtype if available
  \usepackage[]{microtype}
  \UseMicrotypeSet[protrusion]{basicmath} % disable protrusion for tt fonts
}{}
\makeatletter
\@ifundefined{KOMAClassName}{% if non-KOMA class
  \IfFileExists{parskip.sty}{%
    \usepackage{parskip}
  }{% else
    \setlength{\parindent}{0pt}
    \setlength{\parskip}{6pt plus 2pt minus 1pt}}
}{% if KOMA class
  \KOMAoptions{parskip=half}}
\makeatother
\usepackage{xcolor}
\usepackage[margin=1in]{geometry}
\usepackage{color}
\usepackage{fancyvrb}
\newcommand{\VerbBar}{|}
\newcommand{\VERB}{\Verb[commandchars=\\\{\}]}
\DefineVerbatimEnvironment{Highlighting}{Verbatim}{commandchars=\\\{\}}
% Add ',fontsize=\small' for more characters per line
\usepackage{framed}
\definecolor{shadecolor}{RGB}{248,248,248}
\newenvironment{Shaded}{\begin{snugshade}}{\end{snugshade}}
\newcommand{\AlertTok}[1]{\textcolor[rgb]{0.94,0.16,0.16}{#1}}
\newcommand{\AnnotationTok}[1]{\textcolor[rgb]{0.56,0.35,0.01}{\textbf{\textit{#1}}}}
\newcommand{\AttributeTok}[1]{\textcolor[rgb]{0.13,0.29,0.53}{#1}}
\newcommand{\BaseNTok}[1]{\textcolor[rgb]{0.00,0.00,0.81}{#1}}
\newcommand{\BuiltInTok}[1]{#1}
\newcommand{\CharTok}[1]{\textcolor[rgb]{0.31,0.60,0.02}{#1}}
\newcommand{\CommentTok}[1]{\textcolor[rgb]{0.56,0.35,0.01}{\textit{#1}}}
\newcommand{\CommentVarTok}[1]{\textcolor[rgb]{0.56,0.35,0.01}{\textbf{\textit{#1}}}}
\newcommand{\ConstantTok}[1]{\textcolor[rgb]{0.56,0.35,0.01}{#1}}
\newcommand{\ControlFlowTok}[1]{\textcolor[rgb]{0.13,0.29,0.53}{\textbf{#1}}}
\newcommand{\DataTypeTok}[1]{\textcolor[rgb]{0.13,0.29,0.53}{#1}}
\newcommand{\DecValTok}[1]{\textcolor[rgb]{0.00,0.00,0.81}{#1}}
\newcommand{\DocumentationTok}[1]{\textcolor[rgb]{0.56,0.35,0.01}{\textbf{\textit{#1}}}}
\newcommand{\ErrorTok}[1]{\textcolor[rgb]{0.64,0.00,0.00}{\textbf{#1}}}
\newcommand{\ExtensionTok}[1]{#1}
\newcommand{\FloatTok}[1]{\textcolor[rgb]{0.00,0.00,0.81}{#1}}
\newcommand{\FunctionTok}[1]{\textcolor[rgb]{0.13,0.29,0.53}{\textbf{#1}}}
\newcommand{\ImportTok}[1]{#1}
\newcommand{\InformationTok}[1]{\textcolor[rgb]{0.56,0.35,0.01}{\textbf{\textit{#1}}}}
\newcommand{\KeywordTok}[1]{\textcolor[rgb]{0.13,0.29,0.53}{\textbf{#1}}}
\newcommand{\NormalTok}[1]{#1}
\newcommand{\OperatorTok}[1]{\textcolor[rgb]{0.81,0.36,0.00}{\textbf{#1}}}
\newcommand{\OtherTok}[1]{\textcolor[rgb]{0.56,0.35,0.01}{#1}}
\newcommand{\PreprocessorTok}[1]{\textcolor[rgb]{0.56,0.35,0.01}{\textit{#1}}}
\newcommand{\RegionMarkerTok}[1]{#1}
\newcommand{\SpecialCharTok}[1]{\textcolor[rgb]{0.81,0.36,0.00}{\textbf{#1}}}
\newcommand{\SpecialStringTok}[1]{\textcolor[rgb]{0.31,0.60,0.02}{#1}}
\newcommand{\StringTok}[1]{\textcolor[rgb]{0.31,0.60,0.02}{#1}}
\newcommand{\VariableTok}[1]{\textcolor[rgb]{0.00,0.00,0.00}{#1}}
\newcommand{\VerbatimStringTok}[1]{\textcolor[rgb]{0.31,0.60,0.02}{#1}}
\newcommand{\WarningTok}[1]{\textcolor[rgb]{0.56,0.35,0.01}{\textbf{\textit{#1}}}}
\usepackage{graphicx}
\makeatletter
\def\maxwidth{\ifdim\Gin@nat@width>\linewidth\linewidth\else\Gin@nat@width\fi}
\def\maxheight{\ifdim\Gin@nat@height>\textheight\textheight\else\Gin@nat@height\fi}
\makeatother
% Scale images if necessary, so that they will not overflow the page
% margins by default, and it is still possible to overwrite the defaults
% using explicit options in \includegraphics[width, height, ...]{}
\setkeys{Gin}{width=\maxwidth,height=\maxheight,keepaspectratio}
% Set default figure placement to htbp
\makeatletter
\def\fps@figure{htbp}
\makeatother
\setlength{\emergencystretch}{3em} % prevent overfull lines
\providecommand{\tightlist}{%
  \setlength{\itemsep}{0pt}\setlength{\parskip}{0pt}}
\setcounter{secnumdepth}{-\maxdimen} % remove section numbering
\ifLuaTeX
  \usepackage{selnolig}  % disable illegal ligatures
\fi
\IfFileExists{bookmark.sty}{\usepackage{bookmark}}{\usepackage{hyperref}}
\IfFileExists{xurl.sty}{\usepackage{xurl}}{} % add URL line breaks if available
\urlstyle{same}
\hypersetup{
  pdftitle={IECD 2C2024 - Clase 10 - Intervalos de Confianza para la Normal},
  pdfauthor={Gonzalo Barrera},
  hidelinks,
  pdfcreator={LaTeX via pandoc}}

\title{IECD 2C2024 - Clase 10 - Intervalos de Confianza para la Normal}
\author{Gonzalo Barrera}
\date{2024-09-16}

\begin{document}
\maketitle

\hypertarget{actividad}{%
\section{Actividad}\label{actividad}}

Consideremos el siguiente escenario abstracto: se desea estimar por
intervalos los parámetros de una v.a. con distribución normal
\(X \sim N(\mu_0, \sigma^2_0)\). A tal fin, se tienen \(n\)
realizaciones de \(X\): \(\{ x_1, x_2, \dots, x_n \}\).

\begin{enumerate}
\def\labelenumi{\arabic{enumi}.}
\tightlist
\item
  Escriba una función \texttt{ic\_mu\_sigma2\_conocido} que tome tres
  parámetros:
\end{enumerate}

\begin{itemize}
\tightlist
\item
  \texttt{x} (obligatorio), el vector con la muestra
  \(\{ x_1, x_2, \dots, x_n \}\),
\item
  \texttt{sigma2} (obligatorio), el valor de la varianza \(\sigma^2\)
  conocida.
\item
  \texttt{alfa} (por defecto \texttt{0.05}) tal que que
  \(0 < \alpha < 1/2\) y devuelve un intervalo de confianza
  \emph{simétrico en cuantiles} de nivel \$ 1 - \alpha\$ para \(\mu\).
  El intervalo de confianza debe de ser un vector de tipo
  \texttt{"double"} y longitud 2, con la cota inferior del IC en su
  primer posición y la cota superior en la segunda.
\end{itemize}

\begin{Shaded}
\begin{Highlighting}[]
\NormalTok{ic\_mu\_sigma2\_conocido }\OtherTok{\textless{}{-}} \ControlFlowTok{function}\NormalTok{(x, var, }\AttributeTok{alfa=}\FloatTok{0.05}\NormalTok{) \{}
  \CommentTok{\# Su código aquí}
  \CommentTok{\# ic \textless{}{-} c(cota\_inf, cota\_sup)}
  \CommentTok{\# return(ic)}
\NormalTok{\}}
\end{Highlighting}
\end{Shaded}

\begin{enumerate}
\def\labelenumi{\arabic{enumi}.}
\setcounter{enumi}{1}
\item
  Escriba una función \texttt{ic\_mu\_sigma2\_desconocido} siguiendo el
  patrón de del punto 1, pero sin el parámetro \texttt{sigma2}, que
  deberá ser estimado por la propia función.
\item
  Escriba una función \texttt{ic\_sigma2\_mu\_conocido} que tome tres
  parámetros:
\end{enumerate}

\begin{itemize}
\tightlist
\item
  \texttt{x} (obligatorio), el vector con la muestra
  \(\{ x_1, x_2, \dots, x_n \}\),
\item
  \texttt{mu} (obligatorio), el valor de la media \(\mu\) conocida.
\item
  \texttt{alfa} (por defecto \texttt{0.05}) tal que que
  \(0 < \alpha < 1/2\) y devuelve un intervalo de confianza
  \emph{simétrico en cuantiles} de nivel \$ 1 - \alpha\$ para
  \(\sigma^2\). El intervalo de confianza debe de ser un vector de tipo
  \texttt{"double"} y longitud 2.
\end{itemize}

\begin{enumerate}
\def\labelenumi{\arabic{enumi}.}
\setcounter{enumi}{3}
\item
  Escriba una función \texttt{ic\_sigma2\_mu\_desconocido} siguiendo el
  patrón de del punto 3, pero sin el parámetro \texttt{mu}, que deberá
  ser estimado por la propia función.
\item
  Elija unos valores razonables \((\mu_0, \sigma^2_0)\). Fije también
  \texttt{n\_muestra\ \textless{}-\ 10} y
  \texttt{n\_iter\ \textless{}-\ 1000}. Genere una matriz de dimensión
  \((n_{iter}, n_{muestra})\) donde cada fila represente una muestra de
  la v.a. original, con \(n = n_{muestra}\).
\end{enumerate}

\emph{Nota: considere usar los comandos \texttt{matrix} y \texttt{rnorm}
en conjunto para generar los datos en una sola línea de código.}

\begin{enumerate}
\def\labelenumi{\arabic{enumi}.}
\setcounter{enumi}{5}
\item
  Para cada una de las funciones de (1) a (4), estime el parámetro
  relevante a nivel de confianza \(0,95\) para cada una de las
  \texttt{n\_iter} muestras de (5). Guarde los resultados en 4 matrices
  de dimensión \((n_{iter}, 2)\):
  \texttt{est\_mu\_sigma2\_conocido,\ est\_mu...}.
\item
  Calcule el ``nivel de cubrimiento empírico'', es decir, la proporción
  de los IICC generados por cada uno de los métodos, que incluye al
  verdadero parámetro. Qué relación guardan estas proporciones con el
  ``nivel de confianza'' \(1 - \alpha\)? Por qué?
\end{enumerate}

\emph{Nota: Tenga en cuenta que tanto los operadores de comparación
numérica
\texttt{\textless{},\ \textless{}=,\ ==,\ \textgreater{}=,\ \textgreater{}}
como los lógicos \texttt{\textbar{}} (OR) y \texttt{\&} (AND) trabajan
elemento-a-elemento (``element-wise'') sobre los vectores (un objeto
``primitivo'' de \textbf{R}), y \texttt{mean(vec)} devuelve la
proporción de TRUE en \texttt{vec} si
\texttt{typeof(vec)\ ==\ "logical"} (es booleano). Con todo esto, el
cómputo del cubrimiento se reduce a una línea de código.}

\end{document}
